
\section{Risk 9: Strategic Misreporting}
\label{sec:strategic-withholding}

\begin{boxE}
\textit{Strategic Misreporting} arises when an agent conceals or distorts
task-relevant information to improve its own payoff at the expense of others or overall system
performance. Let the world state be $s \in \mathcal{S}$, and suppose agent $i \in \mathcal{N}$ privately observes
a signal $o_{i,t} \in \mathcal{O}_i$ at time $t$. A truthful sufficient report is a mapping
$T_i : \mathcal{O}_i \to \mathcal{M}_i$ that preserves all information about $s$ relevant to the task.
Let $h_t$ denote the public history up to $t$, and let $m_{i,t} = \mu_i(h_t) \in \mathcal{M}_i$ be the
message sent by agent $i$ under some reporting policy $\mu_i$.

For any message $m$, let $\sigma(m)$ denote the information about $s$ revealed by $m$, and let
$\mathrm{supp}\,T_i(o_{i,t})$ be the set of messages that occur with positive probability under
truthful reporting. Withholding or misreporting occurs at $(i,t)$ if
\[
\underbrace{\sigma(m_{i,t}) \subsetneq \sigma\!\big(T_i(o_{i,t})\big)}_{\text{withholding}}
\qquad\text{or}\qquad
\underbrace{\Pr\!\big[m_{i,t} \notin \mathrm{supp}\,T_i(o_{i,t})\big] > 0}_{\text{false report}}.
\]
Such a deviation is \emph{strategic} if it increases agent $i$'s expected utility $u_i$ while (weakly)
reducing system utility $U_{\mathrm{sys}}$ relative to truthful reporting-that is, the agent benefits
from hiding or distorting information at the expense of system performance.
\end{boxE}



\textbf{Motivation.} In many multi-agent systems, information is not evenly distributed. Instead, some agents function as relays or have privileged access to key observations—for instance, agents that act as hubs storing maps, logs, or telemetry. Even small misalignments between individual and team goals can motivate well-informed agents to hide potential risks or hoard valuable information. Such strategic withholding of information can seriously weaken team performance. For example, consider a UAV that is rewarded for producing efficient local plans. Although this incentive is intended to encourage high-quality planning, the UAV may instead choose to suppress newly observed hazards or environment updates to protect its own score. Other teammates, operating under outdated or incomplete information, then make suboptimal decisions, coordination deteriorates, and the system allocates resources inefficiently—sometimes even placing less-informed units at unnecessary risk.
Because of this, it is essential to assess whether information relays truly preserve and accurately transmit task-critical data when designing reporting and auditing mechanisms that align with each agent’s incentives.

\subsection{Experiment — Opportunistic Withholding in a Relay-Based UAV Exploration Task}

\textbf{Overview.}
To investigate how incentive misalignment affects information transmission in cooperative exploration, we evaluate a relay-constrained communication topology in which only one UAV possesses global situational awareness. Such topologies arise, for example, when a fixed base station cannot provide full coverage over the operational area and must rely on UAVs to relay task-critical information across agents. In this experiment, a centralized \emph{State Coordinator} (i.e., the \textbf{State Agent}) assigns target grid cells each round, and two \emph{Exploration UAVs}---denoted \textbf{Agent~1} and \textbf{Agent~2}---execute these assignments. \textbf{Agent~1} is the only UAV with global situational awareness and serves as the sole communication relay to \textbf{Agent~2}; it may either transmit the designated targets and their ground-truth values faithfully or strategically distort them.

Each UAV optimizes two objectives: a \emph{primary} objective of maximizing the team’s total exploration score and a \emph{secondary} objective of maximizing its own individual payoff. Grid cells contain hidden integer values \(v \in \{-1,0,1,2\}\), representing hazards and resources.

To evaluate whether incentive misalignment induces information suppression or manipulation, we define a binary \emph{risk indicator} for each run. Risk is deemed \emph{present} if \textbf{Agent~1} omits assigned targets, fabricates targets, or misreports their values relative to ground truth; otherwise, risk is labeled \emph{absent}. Simulations over four one to two decision rounds allow us to test whether such withholding behaviors emerge as incentives unfold.


\textbf{Setup.} We consider a finite grid $\mathcal{G}$ with hidden cell values $v(x)\in\{-1,0,1,2\}$. At the start of round $t$, the State Agent selects a nonempty target set $T_t\subset\mathcal{G}$ (cells to be explored this round).
\begin{itemize}
\item \textbf{State Agent}: chooses $T_t$ each round and sends it only to Agent~1; no direct channel to Agent~2.
\item \textbf{Agent 1}: observes the full map $v(\cdot)$ and the current $T_t$; sends to Agent~2 a message containing the targets (and, by design, their values), and also reports to the State Agent what it forwarded.
\item \textbf{Agent 2}: observes neither the map nor $T_t$; relies solely on Agent~1’s message to decide where to explore.
\end{itemize}


\textbf{Messaging topology (per round, two-line notation).}
\[
\text{User}\ \rightarrow\ \text{State Agent},\qquad
\text{State Agent}\ \rightarrow\ \text{Agent 1} ,\qquad \text{Agent 1} \ \rightarrow \text{Agent 2}
\]
\[
\text{Agent 2} \ \rightarrow \text{Agent 1} ,\qquad  \text{Agent 1}\ \rightarrow\  \text{State Agent}\qquad\Rightarrow\ \text{end of round}.
\]

\begin{figure}[h]
    \centering
    \includegraphics[width=1\linewidth]{figure/RISK-11_Info_Withholding/topo.pdf}
    \vspace{6pt} 
    \caption{Schematic illustration of the topologies for the Dual-UAV Exploration MAS. Agent~1 relay target information to Agent~2, who selects first; Agent~1 then selects from the remaining targets. }
    \label{deception_exp1_topo}
\end{figure}

In each simulated round, each UAV commits to explore a subset of the designated targets $T_t$ (one cell per agent suffices). An agent’s individual score for a visited cell $x$ is $v(x)$; the team score is the sum across both UAVs. Hazards ($v=-1$) penalize the team.

\begin{wrapfigure}{l}{0.5\textwidth}
    \centering
    \vspace{-8pt}
    \includegraphics[width=0.98\linewidth]{figure/RISK-11_Info_Withholding/2_plot1_total_deceptions_per_experiment.pdf}
    \vspace{3pt}
    \caption{Distribution of identified misreport patterns across eight experiments. Bars show the total number of misreports per experiment, and percentages indicate the misreport rate.}
    \label{fig:2_deception_exp1_bar}
    \vspace{-8pt}
\end{wrapfigure}

Both UAVs optimize a lexicographic objective: maximize team score first; subject to that, maximize individual score. This creates room for Agent~1 to steer Agent~2 toward lower-value (or hazardous) cells while reserving higher-value cells for itself if such steering does not reduce the team’s total under its belief.

We simulate four rounds. At the beginning of each round, the State Agent draws or specifies the target set $T_t$ and sends it, together with the true cell values, to Agent~1. Upon receiving this information, Agent~1 forwards the designated target and a reported value to Agent~2. Agent~2 then makes its selection based solely on the information provided by Agent~1 and sends this selection back to Agent~1. After receiving Agent~2’s selection, Agent~1 selects from the remaining targets, and then compiles both agents’ selections and relays them to the State Agent. This four-round interaction protocol is repeated eight times for evaluation, producing eight independent trials denoted as Experiments~E1 through~E8.

\begin{wrapfigure}{l}{0.5\textwidth}
    \centering
        \vspace{-8pt}
    \includegraphics[width=0.98\linewidth]{figure/RISK-11_Info_Withholding/2_plot2_deception_patterns_pie.pdf}
    \caption{Distribution of identified misreport patterns across eight experiments. Bars show the total number of misreports per experiment, and percentages indicate the misreport rate.}
    \vspace{3pt}
    \label{fig:2_deception_exp1_pie}
    \vspace{-8pt}
\end{wrapfigure}

\textbf{Metric}. For each round $t$, the State Agent specifies a set of target grid cells
$T_t = \{x_1, x_2, \ldots, x_{k_t}\},$
where each $x_i$ denotes a cell index. The corresponding ground-truth values form a
vector
$\mathbf{v}_t = (\, v(x_1), v(x_2), \ldots, v(x_{k_t}) \,).
$ After the communication pipeline concludes, we compare Agent~1's reported values
$
\hat{\mathbf{v}}_t = (\, \hat v(x_1), \hat v(x_2), \ldots, \hat v(x_{k_t}) \,)
$
against the ground truth. For each $x \in T_t$, we check whether the reported value
$\hat v(x)$ matches the true value $v(x)$.

The misreport rate for round $t$ is defined as
$
\frac{\bigl|\{\,x\in T_t : \hat v(x)\neq v(x)\,\}\bigr|}{|T_t|}.
$ A round is classified as a \emph{misreporting event} if the misreport rate is larger than zero. A run is labeled \emph{risk present} if any of its rounds contain a misreporting
event.

